% !TeX root = ../papers/nc-p3.tex
\section{第三長方形の二零点条件が定める曲線の完全分類}
\label{sec:parallel-double-zero-curve}
\begin{theorem}\label{thm:parallel-double-zero-curve}
第三長方形で二つの可変成分 $C,x_{10}$ がともに消える台88の全点が標準Koszul左右AS正則で、
各始点は四つの中心四次式上で階数256の自由加群となる。$A=v,D=2/v,b=2$ とすると、無印の完全不変量は
$X=v^2\in k^*$ である。一位置移動した窓はこの指定始点の曲線上のどの窓とも同型でない。
二位置移動は $X\mapsto4/X$ であり、最小周期は $X=\pm2$ のとき2、他では4である。
反対階数は $2v^2\pm5v+4=0$ の四つのテンソル値で15、他は16である。
\end{theorem}
\begin{proof}[無印不変量と文字直線]
四次交代収縮の値は全曲線で $-16$ である。各始点 $p$ での規格化した収縮
$L_p=W_{p,p+1\mid p+2,p+3}(G/D_4)_{p,p+1\mid p+2,p+3}^t$ の二乗を取り、
後の文字だけを縮約した作用素を $T_p$、前の文字だけを縮約した作用素を $U_{p+1}$ とする。
$T_0$ の跡は $-157/16$ で、$T_1$ の跡は $(v^2+4/v^2+316)/16$ である。
後者だけで同型候補は $v'\in\{v,-v,2/v,-2/v\}$ に限られる。前者の二乗跡は
\[
 \operatorname{tr}(T_0(v)^2)-\operatorname{tr}(T_0(2/v)^2)
 =\tfrac{55}{1024}(v^2-4/v^2).
\]
従って元から二乗が等しい場合だけが残る。$v\mapsto-v$ は四位置の文字平行移動量 $(3,2,2,3)$ と対角倍率で実現する。
二位置移動と $v\mapsto2/v$ の合成は全位置で文字0,1を交換した後の対角倍率で実現する。
実際、共通台の語の係数比を取り、位置別倍率の指数行列の左整数核の全基底について指標値が1である。

一位置移動については、$T_p,U_p$ の四組の同時固有値を使える。
いずれの位置でも、異なる二文字での二つの値の差の分子は互いに素なので四同時固有直線を全点で無印復元できる。
元の台の任意の位置別文字置換で一位置移動後の台になるものはない。
これは初めの三位置の $24^3$ 通りを取り、総色零から第四位置でなければならない文字を決める有限検査である。
従って一位置移動でこの曲線に戻ることはない。
反対四ブロックの行列式は $-(2v^2-5v+4)/v,(2v^2+5v+4)/v,-1,-1$ となり、四根は相異なり単純である。
\end{proof}

\begin{proof}[関数体の証明を外れた全根の処理]
前節でこの曲線の点 $v=1$ が正則であるため、曲線の幾何学的一般点も正則で、
各始点・次数のHilbert関数は全特殊点で標準値以上である。有理関数体 $\mathbb Q(v)$ 上で
五次までの経路基底と四次中心を計算する。文字0では三つの自由列で係数 $(1,3,5),(2,7,11)$ の二式、
文字1では二自由列の基底を取り、これらの中心商を四始点で長さ13まで計算すると
$(1+z+z^2+z^3)^4$ の係数を得る。中心四式で割った後は、基底語に一矢を付けて末尾二次関係を課し、
長さ4で始点の中心四式を加える。さらに長い語での中心式は始点に移せるので同じ商である。

この計算が使える開集合を厳密に記録する。経路基底・中心の核・中心商を選ぶ全行列で分母を払い、
行の内容を割る際にはその非定数因子を記録する。主元および後退消去後の各主元の非定数因子も記録する。
一列の消去で二係数の最大公約数を割るのは主元が可逆な局所化での変形である。
従って記録した因子が全て非零なら、全ての基底・次元・中心式がそのまま特殊化する。
標準Hilbert下界と自由加群による上界が一致し、ここでは全次数で正則である。

記録される既約因子は1885個である。ただし正則性は $v\mapsto-v,2/v$ の両方で保たれるため、
その軌道全体が例外集合に入る根だけを調べれば十分である。零を除くと40個の既約多項式、計64根になる。
各多項式 $f$ について先頭係数が可逆な素数 $l$ と非零の一次剰余根を選ぶ。
五次までの経路次元はその剰余根でも標準で、四次中心の文字別次元は $(3,2,0,0)$ である。
剰余体上でその核から四式を取り、全四始点で同じ長さ13までの中心商次元を検算する。

この剰余計算は $f$ の標数零の根での証明になる。実際、数体でその剰余根の上の素点を取り、
経路空間の有限次数の行列を比較する。長さ5まで、曲線の一般点からの下界と剰余体からの上界が等しいため、
各次数は素点で自由である。この基底で四次中心の式を立てると、標数零の核次元も曲線の一般値以上であり、
剰余体の核次元と等しくなる。よって中心核はその素点で直和因子で、指定の剰余四式を中心式として持ち上げられる。
中心商には剰余体からの次元上界が移り、元の全次数のHilbert下界と比較すれば自由性と正則性が従う。
同じ既約多項式の全ての共役根に適用できる。
\path{computations/verify_parallel_double_zero_curve.py} と
\path{computations/verify_parallel_curve_residues.py} が開集合と全残余根を再生成する。
\end{proof}
